\section{The Critic and the Ideologue}

\begin{quotex}
We will lower the screens and drop the veils on the eye of the critic [one who uses logical proof] and the one who simply follows a creed [ideologue]. Both of them are among the forms in which God gives tajalli [his Self-manifestation], but He has commanded the veil for us in order that different levels of excellence of the predisposition of forms might appear. \flright{\textsc{Ibn Arabi}, \emph{Bezels of Wisdom}}

\end{quotex}
Recently, there was a comment on the \textit{Five Elements of Manifestation}\footnote{\url{https://www.gornahoor.net/?p=14558}}:

\begin{quotex}
While it is commendable to do amateur philosophy, but to include Eliphas Levi with Indian philosophy is unpardonable.

\end{quotex}
First of all, if your talking donkey tells you the truth, then it is still true, and a professional intellectual may talk nonsense even with a million followers. Levi was just summarizing \textbf{Paracelsus}, who has a better reputation, but it is easier to be snarky than to be informed. Nor is there any mention of ``Indian" philosophy, because metaphysics transcends nationality. Either the five elements are qualitatively accurate everywhere or they are not so anywhere. We are not ``doing philosophy" but describing reality. You see it or you don't; the critic is veiled with illusion, but he is doing the best he can. God knows best what degree of self-manifestation is appropriate to each person.

A ``professional" philosopher will tell you all about Plato, his various periods, his teachings, and so on. He may even mention that Plato taught that we live in a world of shadows and few bother to look for the Sun that casts the shadows. A fortiori, he does not even know how to escape the shadow world to the enlightened world.

The philosopher will read that there is a form of knowing, called ``Intuition", that is superior to what the rational mind can achieve. This is the case in metaphysics from India, to Greece, to Egypt, to Arabia, to the Scholastics. Yet the veil blocks his understanding so he ignores it. Intuition is akin to ``seeing", not to ``thinking", and certainly not to ``arguing".

\paragraph{The Culinary Arts}
Plato complained that the purpose of the culinary arts was to make unpalatable or unhealthy foods taste delicious. Analogously, the professionals try to make life in the cave palatable. The world is replete with distractions that make it possible to get through the day. (I leave it to the reader to make a personal list of such distraction.)

What does the one who has seen the Sun then do? He has become a gnostic, ``one who knows". Does he leave the shadows behind or does he go back to find some others, always a small minority, who would like to see the Sun? Perhaps you have an opinion. Keep in mind Aristotle's sobering claim that there are those—perhaps most—who are simply unteachable.

\paragraph{Professional Opinions}
I don't offer anything new, since I have referenced dozens of the best minds; I hope some readers have been motivated to return to the original sources. Perhaps some of the illustrations are novel; a diligent study may be able to repeat texts, but if he can't illustrate concrete examples, his understanding may be incomplete.

Physics studies phenomena and tries to derive principles the relate them to each other. Metaphysics starts with principles and understands phenomena through them. I will give examples from thinkers who have a much more sterling reputation than I. \textbf{Sayyid Nasr} explains it this way (using Islamic terminology):

\begin{quotex}
The term ``ta'wil," which plays a cardinal role in Shi'ism as well as Sufism, means literally to return to the origin of a thing. It means to penetrate the external aspect of any reality, whether it be sacred scripture or phenomena of nature, to its inner essence, to go from the phenomenon to the noumenon.

\end{quotex}
Hence, in that article, I started with the inner essence of matter, in its five forms, and illustrated how they correspond to various phenomena. I showed how the elemental beings are experienced as various psychic states. In this regard, I am simply developing ideas from Nasr:

\begin{quotex}
In the religious cosmos of the traditional Muslim, which is filled with material, psychic, and spiritual creatures of God, the jinn play their own particular role. By the elite they are taken for what they are, namely, psychic forces of the intermediate world of both a beneficent and an evil nature: On the popular level, the jinn appear as concrete physical creatures of different shapes and forms against which men seek the aid of the Spirit, often by chanting verses of the Quran.

On the theological and metaphysical level of Islam, the order of the jinn becomes under stood as a necessary element in the hierarchy of existence, an element which relates the physical world to higher orders of reality.

\end{quotex}
So apart from the jinn, there are other semi-material, psychic, and spiritual creatures. Obviously, the semi-material creature—in the sense that they do not incarnate with all the elemental qualities—manifest as various psychic forces. These forces are known to Indian philosophy, Islamic metaphysics, and, yes, even to Eliphas Levi.

\paragraph{Music}
The ethereal nature of music was noted by the professional psychologist, \textbf{Emma Jung}:

\begin{quotex}
For music can be understood as an objectification of the spirit; it does not express knowledge in the usual logical, intellectual sense, nor does it shape matter; instead, it gives sensuous representation to our deepest associations and most immutable laws. In this sense, music is spirit, spirit leading into obscure distances beyond the reach of consciousness; its content can hardly be grasped with words — but strange to say, more easily with numbers — although simultaneously, and before all else, with feeling and sensation. Apparently paradoxical facts like these show that music admits us to the depths where spirit and nature are still one — or have again become one. 

\end{quotex}
You can choose to ``do" philosophy in the ``usual logical, intellectual sense", just like a professional, or you can follow music into ``the depths where spirit and nature are still one — or have again become one."

\paragraph{Jinns}
The related article, \textit{Passing Through the Realms of the Jinns}\footnote{\url{https://www.gornahoor.net/?p=14515}}, illustrates several of these points. There is no essential distinction between that description and that of the professional Nasr:

\begin{quotex}
Man was made of clay. The jinn in Islamic doctrines are that group of creatures which was made of fire rather than earth.

\end{quotex}
Actually, man was made of wet clay: earth and water, as we wrote. Nasr continues:

\begin{quotex}
[Jinns] possess a volatile and ``unfixed" outer form and so can take on many shapes. This means that they are essentially creatures of the psychic rather than the physical world and that they can appear to man in different forms and shapes. 

\end{quotex}
As we pointed out, the fire element has essence but no particular form, more that you get from Nasr. Moreover, that article explains in detail why the jinn would appear in a particular form to a specific person.

\paragraph{The Gnostic}
Ibn Arabi describes the one whose understanding passes beyond reasoning:

\begin{quotex}
When he is alone with himself after God's self-manifestation [tajalli], he is confused about what he saw. If he is the slave of a Lord, his intellect returns to him, and if he is the slave of reasoning, God gives him back judgement. This is only as long as he is veiled in the dimension of this world from his dimension in the Next World. The gnostics appear here as if they possessed the form of this world by virtue of the fact that its principles operate on them. However, God has transformed them inwardly to the dimension of the Next World. That must be the case. \emph{They are unknown by form except to whoever has had his inner eye unveiled by God so that he perceives}. 

\end{quotex}
In other words, you won't be able to recognize the gnostic through his outer appearance alone. I know that some authors like to pose sideways with a vacant stare into the unknown. That counts for little.



\flrightit{Posted on 2021-07-18 by Cologero }

\begin{center}* * *\end{center}

\begin{footnotesize}\begin{sffamily}



\texttt{Cologero on 2021-07-19 at 06:38 said: }

For more on how ``professional" philosophers go wrong, check out \textit{The particle collection that fancied itself a physicist}\footnote{\url{https://edwardfeser.blogspot.com/2020/08/the-particle-collection-that-fancied.html?m=1}} by Edward Feser:

\begin{quotex}
It can be charming when a child pretends that he is a cloud, or a boulder, or a lion, or even – I suppose – a collection of particles. It's considerably less charming when a grown man does it, and when he is a grown man with a Ph.D. and a tenured position at Columbia, it's downright embarrassing. But you need only turn on the news to see that otherwise intelligent people believing ever more ridiculous things on the basis of ever more convoluted sophistries is the story of our age. There is a crucial but widely overlooked lesson here. When your basic assumptions are unsound, greater intelligence by no means guarantees that you will come to see this. On the contrary, sometimes you will end up only more hardened in error than a less intelligent person would be, because you will be able to come up with subtler fallacies and cleverer self-deceptions. 

\end{quotex}

\end{sffamily}\end{footnotesize}
